% Source: https://artofproblemsolving.com/community/c3h2439042_question_about_quadratics
% Stable identity: aops-topic-2439042-q01
\textbf{(a)} The number $r$ is a root if substituting $x=r$ makes the equation true: $r^2-r-6=0$. (Equivalently, $r$ is a zero of the polynomial $x^2-x-6$.)

\textbf{(b)} Factor: $x^2-x-6=(x-3)(x+2)$. A product is $0$ exactly when one of the factors is $0$, so the roots are $\boxed{x=3\text{ and }x=-2}$. Check: $9-3-6=0$ and $4+2-6=0$.

\textbf{(c)} A point of the graph lies on the $x$-axis exactly when $y=0$, i.e. when $x$ is a root. So the roots are the $x$-coordinates of the points where the parabola crosses the $x$-axis, $(-2,0)$ and $(3,0)$.

\begin{asy}
usepackage("amsmath");
size(5cm);
defaultpen(fontsize(10pt));
import graph;
real f(real x){return x^2-x-6;}
draw((-3.6,0)--(4.6,0),Arrow(6)); draw((0,-7.2)--(0,5),Arrow(6));
label("$x$",(4.6,0),E); label("$y$",(0,5),N);
draw(graph(f,-3.1,4.1),blue+linewidth(0.9));
dot((-2,0)); dot((3,0)); label("$(-2,0)$",(-2,0),SW); label("$(3,0)$",(3,0),SE);
label("$y=x^2-x-6$",(3.4,f(3.4)),E,blue);
\end{asy}
